\begin{exercise}{94}
Let $\mathcal C$ be an ordinary category and $S$ a set of morphisms. Show that
$\mathcal C\to h\mathcal C[S^{-1}]$ is the initial functor $\mathcal C\to\mathcal D$
(up to natural isomorphism) between ordinary categories sending $S$ to isomorphisms.

\begin{proof}
Let $\mathcal O$ be the ordinary category obtained from $\mathcal C$ by adjoining
formal inverses of the arrows in $S$, subject to the composition relations of
$\mathcal C$ and the two inverse relations for each $s\in S$.
An $S$-inverting functor extends uniquely to $\mathcal O$ by sending each
formal inverse to the inverse of its image.
Natural transformations extend uniquely as well: their naturality squares for
the inverse arrows follow by multiplying the squares for $s$ by the inverses.
Thus restriction identifies $\mathrm{Fun}(\mathcal O,\mathcal E)$ with the
full subcategory $\mathrm{Fun}_S(\mathcal C,\mathcal E)$ of $S$-inverting functors.

We now verify that $h\mathcal C[S^{-1}]$ has this universal property.
Regard $\mathcal C$ as its nerve and put $L=\mathcal C[S^{-1}]$.
For every simplicial set $X$ and ordinary category $\mathcal E$, there is a
natural isomorphism
\[
 \mathrm{Fun}(X,\mathrm N\mathcal E)
 \cong\mathrm N\bigl(\mathrm{Fun}(hX,\mathcal E)\bigr).
\]
Indeed, by Proposition~1.2.18, Exercise~\ref{ex:28} and Lemma~1.2.7,
an $n$-simplex on either side is a functor $[n]\times hX\to\mathcal E$;
these identifications commute with the simplicial operators.
Under this isomorphism and $h\mathrm N\mathcal C\cong\mathcal C$
(Corollary~1.2.14), Remark~2.4.3 gives a Joyal equivalence
\[
 \mathrm N\bigl(\mathrm{Fun}(hL,\mathcal E)\bigr)
 \longrightarrow\mathrm N\bigl(\mathrm{Fun}_S(\mathcal C,\mathcal E)\bigr)
\]
induced by restriction along $\mathcal C\to hL$.
By Exercise~\ref{ex:89} and Corollary~1.2.14, the corresponding ordinary
restriction functor is an equivalence of categories.
For every $S$-inverting functor this gives a factorization through $hL$,
unique up to natural isomorphism, as required.
In particular, $hL$ and the word category $\mathcal O$ are equivalent under
$\mathcal C$.
\end{proof}
\end{exercise}
